%% start of file `template.tex'.
%% Copyright 2006-2013 Xavier Danaux (xdanaux@gmail.com).
%
% This work may be distributed and/or modified under the
% conditions of the LaTeX Project Public License version 1.3c,
% available at http://www.latex-project.org/lppl/.


\documentclass[11pt,a4paper,sans]{moderncv}        % possible options include font size ('10pt', '11pt' and '12pt'), paper size ('a4paper', 'letterpaper', 'a5paper', 'legalpaper', 'executivepaper' and 'landscape') and font family ('sans' and 'roman')

% moderncv themes
\moderncvstyle{casual}                             % style options are 'casual' (default), 'classic', 'oldstyle' and 'banking'
\moderncvcolor{blue}                               % color options 'blue' (default), 'orange', 'green', 'red', 'purple', 'grey' and 'black'
%\renewcommand{\familydefault}{\sfdefault}         % to set the default font; use '\sfdefault' for the default sans serif font, '\rmdefault' for the default roman one, or any tex font name
\nopagenumbers{}                                  % uncomment to suppress automatic page numbering for CVs longer than one page

% character encoding
\usepackage[utf8]{inputenc}                       % if you are not using xelatex ou lualatex, replace by the encoding you are using
%\usepackage{CJKutf8}                              % if you need to use CJK to typeset your resume in Chinese, Japanese or Korean

% adjust the page margins
\usepackage[scale=0.8]{geometry}
\setlength{\hintscolumnwidth}{3.2cm}                % if you want to change the width of the column with the dates
%\setlength{\makecvtitlenamewidth}{10cm}           % for the 'classic' style, if you want to force the width allocated to your name and avoid line breaks. be careful though, the length is normally calculated to avoid any overlap with your personal info; use this at your own typographical risks...

\makeatletter
\newcommand*{\skypesocialsymbol}{\includegraphics[height=.7\baselineskip]{skype_gray}}
\collectionadd[skype]{socials}{ daniel.fulop@hotmail.com}
\makeatother

% personal data
\name{Daniel}{F{\"u}l{\"o}p}
\email{nearnaln@gmail.com}                               % optional, remove / comment the line if not wanted
\phone[mobile]{+41~77~981~0697}                   % optional, remove / comment the line if not wanted; the optional "type" of the phone can be "mobile" (default), "fixed" or "fax"
\photo[96pt][0.4pt]{picture}                       % optional, remove / comment the line if not wanted; '64pt' is the height the picture must be resized to, 0.4pt is the thickness of the frame around it (put it to 0pt for no frame) and 'picture' is the name of the picture file

% to show numerical labels in the bibliography (default is to show no labels); only useful if you make citations in your resume
%\makeatletter
%\renewcommand*{\bibliographyitemlabel}{\@biblabel{\arabic{enumiv}}}
%\makeatother
%\renewcommand*{\bibliographyitemlabel}{[\arabic{enumiv}]}% CONSIDER REPLACING THE ABOVE BY THIS

% bibliography with mutiple entries
%\usepackage{multibib}
%\newcites{book,misc}{{Books},{Others}}
%----------------------------------------------------------------------------------
%            content
%----------------------------------------------------------------------------------
\begin{document}
%\begin{CJK*}{UTF8}{gbsn}                          % to typeset your resume in Chinese using CJK
%-----       resume       ---------------------------------------------------------
\makecvtitle
\setlength{\parskip}{0.5em}
\section{Pers{\"o}nliche Daten}
\cvitem{Name}{Daniel F{\"u}l{\"o}p}{}
\cvitem{Staatsangeh{\"o}rigkeit}{ungarisch}{}
\cvitem{E-Mail}{nearnaln@gmail.com -- bevorzugt bei der ersten Kontaktaufnahme}{}
\cvitem{Telefon}{+41~77~981~0697 -- bevorzugte Anrufzeit: zwischen 9 und 10 Uhr morgens}{}
\cvitem{Skype}{daniel.fulop@hotmail.com}
\cvitem{Wohnort}{Budapest (bereit in die Schweiz umzuziehen)}{}

\section{Berufserfahrung}
\cventry{Juli 2015 -- \\ Mai 2016}{Software-Entwickler}{evosoft Hungary Kft}{}{}{Entwicklung einer IoT-Plattform mit Microservices mithilfe des Spring Clouds.\newline{}%
Java (Spring Cloud), Microservices, Cloud Foundry, OAuth2, Gradle
\newline{}
Aufgaben:
\begin{itemize}
  \item Zusammenarbeit mit fernen Teams der Siemens AG aus Deutschland
  \item Verwendung des OAuth2-Protokolls und seiner Implementation im Spring Framework.
  \item Installation von Java-Anwendungen auf CloudFoundry, Fehlersuche und -behebung im Cloud.
  \item Umstellung auf einen anderen CloudFoundry-Provider unter Zusammenarbeit mit lokalen und fernen Entwicklungsteams.
  \item Teilnahme an der Ausarbeitung von technologischen und architektonischen Pl{\"a}nen mit fernen Teams.
\end{itemize}}

\cventry{Dezember 2014 -- \\Juli 2015}{Software-Entwickler}{Nextent Informatika Zrt. und Erste Bank Hungary Zrt}{}{}{Unterst{\"u}tzung des Softwaretestens durch Entwicklung von Testtools.\newline{}%
Python, .NET (WPF), Oracle Database 
\newline{}
Aufgaben:
\begin{itemize}
  \item Implementation einer VBA-Anwendung in Excel f{\"u}r das Pr{\"u}fen von Finanzberichten.
  \item Entwicklung von Python-Tools f{\"u}r das Testen und Verarbeitung von SQL-Skripten.
  \item Programmierung einer .NET WPF-Anwendung f{\"u}r Daten{\"u}bertragung zwischen Datenbankschemata.
  \item Entwicklung eines Subversion-Clients um die spezifischen Aufgaben des Entwicklerteams zu unterst{\"u}tzen.
\end{itemize}}

\cventry{April 2013 -- \\Mai 2014}{Software-Entwickler}{IP Systems Kft}{}{}{Entwicklung einer Plattform f{\"u}r Erdgastransport und -handel.\newline{}%
Java (Spring Framework), Oracle Database, Subversion, Maven
\newline{}
Aufgaben:
\begin{itemize}
  \item Fehlersuche und -behebung unter SLA.
  \item Demonstration des Entwicklungsfortschrittes alle zwei Wochen an Kunden.
  \item Zusammenarbeit mit Teams von unterschiedlichen Produkten um technologische Erfahrungen auszutauschen.
\end{itemize}}

\section{Ausbildung}
\cventry{2011 -- 2016}{Informatik BSc}{Technische und Wirtschaftswissenschaftliche Universit{\"a}t Budapest}{}{}{Fachrichtung: Softwareentwicklung}


\section{Kompetenzen}
\cvitem{Programmiersprachen}{\textbf{Sehr gut:} Java, SQL\newline{}\textbf{Gut:} Python, JavaScript\newline{}\textbf{Anf{\"a}nger:} C\#}{}
\cvitem{Frameworks}{\textbf{Sehr gut:} Spring Framework\newline{}\textbf{Anf{\"a}nger:} .NET (MVC, Web API), WPF}{}
\cvitem{Architekturen}{REST, Microservices}{}
\cvitem{Plattformen}{Cloud Foundry}{}
\cvitem{Methoden}{Scrum, TDD, Paarprogrammierung}{}

\section{Sprachkenntnisse}
\cvlistitem{Deutsch -- verhandlungssicher, C1}
\cvlistitem{Englisch -- verhandlungssicher, C1}

\section{Referenzen}
\cvitem{}{Auf Anfrage}{}

\end{document}

%% end of file `template.tex'.
